\documentclass[10pt,a4paper]{amsart}
\usepackage[margin=19mm]{geometry}
\usepackage{amsmath,amssymb,mathtools}
\usepackage[scr=rsfs,cal=euler]{mathalfa}
\usepackage{xfrac}
\usepackage[unicode,hidelinks,bookmarks=false]{hyperref}
\newcommand{\ie}{i.e.\ }
\newcommand{\cf}{cf.\ }
\newcommand{\resp}{respectively\ }
\newcommand{\Subsection}{\subsection}
\providecommand{\nobreakdash}{\nobreak\hbox{-}}
\setlength{\emergencystretch}{2em}
\multlinegap0pt
\usepackage{mdframed}
\usepackage{mdframed}
\usepackage{mdframed}
\usepackage{mdframed}
\usepackage{amssymb}
\usepackage{mathtools}
\usepackage[shortlabels]{enumitem}
\usepackage{xcolor}
\usepackage{csquotes}
\usepackage{mdframed}
\usepackage{tikz}
\newtheorem{lemma}{Lemma}
\title{From Computational Certification to Exact Coordinates:\\ Heilbronn's Triangle Problem on the Unit Square\\ Using Mixed-Integer Optimization}
\author{Nathan Sudermann-Merx}
\date{Source: arXiv:2603.11107v2 (CC BY 4.0)}
\begin{document}
\maketitle

\noindent\emph{Source and licence.} From Computational Certification to Exact Coordinates:\\ Heilbronn's Triangle Problem on the Unit Square\\ Using Mixed-Integer Optimization. Version arXiv:2603.11107v2 (CC BY 4.0), distributed by arXiv under the Creative Commons Attribution 4.0 International licence, http://creativecommons.org/licenses/by/4.0/, as recorded in the arXiv licence metadata for this exact version and shown on the abstract page https://arxiv.org/abs/2603.11107v2. Full licence notice: \texttt{licenses/arxiv-2603.11107-CC-BY-4.0.txt}.\par\smallskip

\noindent\textit{Editorial scope.} Counted unit: the article's lemma lem:min\_parallelogram (source0.tex lines 407--409). The article's complete original proof of it is delivered with it (source0.tex lines 411--437, one \texttt{proof} environment of 27 lines). No mathematics is written, repaired or re-derived: the statement and the proof are the article's own lines, in the article's own notation, and no retained line is substituted or re-flowed.  No further source-local context is needed: every cross-reference of every retained line resolves inside the two delivered spans.  Nothing was omitted from any retained span, so the package carries no omission record.  No retained line contains a citation, so the wrapper carries no bibliography and no bibliography entry.  No retained line contains a figure environment, an image-inclusion command or drawing code, so no image is delivered and the package declares no asset dependency.  Editorial formatting only, in addition: an AMS \texttt{amsart} preamble that restates the article's own declarations and macros used by the retained lines, namely the environments \texttt{lemma} on separate counters, together with the standard packages those lines need.  The retained environments print in this document as Lemma 1 in order of appearance, whereas the article prints the same items with the numbering of its own full document.  The complete native source of the article as distributed in the arXiv e-print for this exact version is bundled unchanged as \texttt{sources/196193/source0.tex}.
\begin{lemma}[Minimal covering parallelogram]\label{lem:min_parallelogram}
The unit square $S=[0,1]^2$ is a minimum-area covering parallelogram of $K$.
\end{lemma}

\begin{proof}
Assume for contradiction that there exists a parallelogram $S' \supseteq K$
with $\operatorname{area}(S') < 1$.

Since every parallelogram is an affine image of the unit square,
there exists an affine bijection $T: S' \to S$
with positive determinant satisfying
\[
|\det T| = \frac{\operatorname{area}(S)}{\operatorname{area}(S')}
= \frac{1}{\operatorname{area}(S')} > 1.
\]

Because $K \subseteq S'$, we have $T(K) \subseteq S$, hence $T(P) \subseteq S$.
For every triangle $\tau$ formed by points of $P$,
its area scales by
\[
\operatorname{area}(T(\tau))
= |\det T| \operatorname{area}(\tau).
\]
Thus
\[
\Delta(T(P))
= |\det T| \Delta(P)
> \Delta(P),
\]
contradicting the optimality of $P$.
\end{proof}

\end{document}
